\documentclass[11pt,a4paper]{article}

\input{../../sablony/pripravy/priprava.sty}

\nastavPripravu{4. ročník}{Kombinatorika}{34}{Variace I -- základy}

\begin{document}

\blok{Motivace a návaznost}

Ze sedmi uchazečů obsadíme tři různá místa na stupních vítězů. Na rozdíl od permutací nepoužijeme všechny uchazeče.

\poznamka{Předpoklady: permutace bez opakování, faktoriál, pravidlo součinu. Cíl: pochopit uspořádaný výběr jen části prvků a odvodit vzorec variací bez opakování.}

\blok{Výklad a odvození}

\begin{definiceBox}
\textbf{Variace bez opakování} délky $k$ z $n$ různých prvků je uspořádaný výběr $k$ různých prvků. Pro celá $0\le k\le n$ platí
\[
V(n,k)=n(n-1)\cdots(n-k+1)=\frac{n!}{(n-k)!}.
\]
\end{definiceBox}
U permutace řadíme všech $n$ objektů; u variace jen $k$ z nich a rozlišujeme pořadí nebo role. Pro $k=n$ dostaneme permutaci. Je-li $k=0$, existuje jediný prázdný výběr.

\blok{Společně řešené příklady}

\begin{prikladReseniBox}[title={Příklad 1}]
Osm běžců obsadí první tři místa bez dělených umístění. Kolik výsledků existuje?
\tcblower
$V(8,3)=8\cdot7\cdot6=336$.
\end{prikladReseniBox}



\begin{prikladReseniBox}[title={Příklad 2}]
Ze šesti žáků vybereme předsedu a místopředsedu. Kolik obsazení funkcí?
\tcblower
$V(6,2)=6\cdot5=30$. Funkce jsou různé.
\end{prikladReseniBox}



\begin{prikladReseniBox}[title={Příklad 3}]
Kolik pětimístných řetězců bez opakování lze vytvořit ze sedmi různých písmen?
\tcblower
$V(7,5)=7\cdot6\cdot5\cdot4\cdot3=2520$.
\end{prikladReseniBox}



\blok{Samostatná práce}

\begin{zadaniReseniBox}[title={Samostatná práce 1}]
Deset lidí soutěží o zlato a stříbro. Kolik dvojic vítězů?
\tcblower
$10\cdot9=90$.
\end{zadaniReseniBox}



\begin{zadaniReseniBox}[title={Samostatná práce 2}]
Kolik třípísmenných hesel bez opakování vytvoříme ze čtyř různých písmen?
\tcblower
$4\cdot3\cdot2=24$.
\end{zadaniReseniBox}



\begin{zadaniReseniBox}[title={Samostatná práce 3}]
Vysvětli rozdíl mezi počtem pořadí šesti lidí a obsazením tří rolí ze šesti lidí.
\tcblower
Permutace: $6!=720$; variace: $6\cdot5\cdot4=120$.
\end{zadaniReseniBox}



\blok{Rozšíření}

\begin{zadaniReseniBox}[title={Rozšíření 1}]
Z $12$ finalistů vybereme první čtyři místa v pořadí. Jaký je výsledek ve faktoriálovém tvaru?
\tcblower
$V(12,4)=12!/8!=12\cdot11\cdot10\cdot9=11880$.
\end{zadaniReseniBox}



\blok{Domácí úkol}

\poznamka{Doporučená práce doma: přibližně 15--30 minut. Úlohy jsou seřazeny od základních k náročnějším; učitel může poslední podúlohu označit jako volitelnou.}

\begin{zadaniReseniBox}[title={Domácí úkol}]

\begin{enumerate}[label=\alph*),leftmargin=*,itemsep=2pt]

\item Kolik různých pořadí prvních tří soutěžících může nastat mezi devíti soutěžícími?
\item Ze sedmi studentů obsadíme tři různé funkce. Kolik možností?
\item Kolik různých čtyřznakových kódů bez opakování vytvoříme z pěti různých písmen?
\item Vysvětli, kdy se ze vzorce $V(n,k)$ stává vzorec pro permutace.
\item Urči přesnou hodnotu $V(9,4)$.
\item Vypočti $V(7,0)$ a $V(7,7)$; vysvětli, co oba výsledky znamenají.

\end{enumerate}

\tcblower

\begin{enumerate}[label=\alph*),leftmargin=*,itemsep=2pt]

\item $9\cdot8\cdot7=504$.
\item $7\cdot6\cdot5=210$.
\item $5\cdot4\cdot3\cdot2=120$.
\item Při $k=n$ máme $V(n,n)=n!=P(n)$.
\item $9\cdot8\cdot7\cdot6=3024$.
\item $1$ (prázdný výběr) a $7!=5040$ (permutace všech sedmi).

\end{enumerate}

\end{zadaniReseniBox}

\poznamka{Orientační práce v hodině: 0--12 min kontrast s permutací; 12--25 min odvození; 25--40 min procvičení; 40--45 min kontrola.}

\end{document}
